% !PRINTED 58
\weberpage{58}

%<<CONT
tion wäre in der ganzen $u$-Ebene endlich und müßte also nach
§\,15, 8 eine Konstante sein.

Ist $\varphi(u)$ eine doppeltperiodische Funktion mit den Perioden
$\omega_1$, $\omega_2$, so ist das über die Begrenzung des Periodenparallelogramms
 genommene Integral [wegen (1)]

\[
\int d\log\varphi(u) = \int\limits_{u_0}^{u_0+\omega_1} [d\log\varphi(u) - d\log\varphi(u+\omega_2)]
\]
\[
{}-\int\limits_{u_0}^{u_0+\omega_2} [d\log\varphi(u) - d\log\varphi(u+\omega_1)] = 0\,.
\]

Daraus folgt nach §\,15, 6. der Satz:

2. \gesperrt{Eine doppeltperiodische Funktion wird in einem
Periodenparallelogramm in ebenso vielen Punkten Null
wie unendlich.} Ist $m$ diese Zahl, so heißt $m$ die Ordnung
der doppeltperiodischen Funktion. Nach 1. gibt es keine
doppeltperiodischen Funktionen von der Ordnung Null.

Ersetzt man in dem Beweise die Funktion $\varphi(u)$ durch $\varphi(u)-c$,
so folgt, daß eine doppeltperiodische Funktion $m$ter Ordnung
jeden beliebigen Wert $c$ in gleich vielen Punkten annimmt.

Wenn wir das Integral der Formel in 7., §\,15 über die Begrenzung
 des Periodenparallelogramms nehmen, so ergibt sich:

\[
\int u\, d\log\varphi(u) = \int\limits_{u_0}^{u_0+\omega_1} [u\,d\log\varphi(u) - (u+\omega_2)\,d\log\varphi(u+\omega_2)]
\]
\[
{}-\int\limits_{u_0}^{u_0+\omega_2} [u\,d\log\varphi(u) - (u+\omega_1)\,d\log\varphi(u+\omega_1)]
\]
\[
{}=\omega_1\int\limits_{u_0}^{u_0+\omega_2} d\log\varphi(u) - \omega_2\int\limits_{u_0}^{u_0+\omega_1} d\log\varphi(u)\,.
\]

Nun ist wegen der Mehrdeutigkeit des Logarithmus

\[
\int\limits_{u_0}^{u_0+\omega_2} d\log\varphi(u) = \log\varphi(u+\omega_2) - \log\varphi(u) = 2\pi i\, m_1
\]
\[
\int\limits_{u_0}^{u_0+\omega_1} d\log\varphi(u) = \log\varphi(u+\omega_1) - \log\varphi(u) = -2\pi i\, m_2\,,
\]

worin $m_1$ und $m_2$ unbestimmte ganze Zahlen sind.

Wenn daher die Funktion $\varphi(u)$ in den Punkten $\alpha_1$, $\alpha_2$, $\dots$, $\alpha_m$
unendlich, in den Punkten $\beta_1$, $\beta_2$, $\dots$, $\beta_m$ Null wird, so ist nach
§\,15, 7.:

\[ \Sigma\alpha_i \equiv \Sigma\beta_i \ (\mathrm{mod}\ \omega_1,\ \omega_2)\,. \]
